\section{Outcomes}
\label{sec:downstream-outcomes}

The administrative results in Sections~\ref{sec:electorate-size} and~\ref{sec:decomposition} establish two facts that organize the remainder of the paper. First, strict biometric registration contracts the registered electorate at adoption, and the contraction is concentrated among voters recorded as low education. Second, part of the compositional shift is not exclusion but re-labeling: voters who pass through recadastramento update stale schooling records, so the high-education count rises even as the total registry falls. These two channels point in the same direction for the object that politicians and local administrators actually observe. After BVR, the active electorate visible in official records is smaller, more selected toward more educated voters, and, in either interpretation, easier to read against the resident population. The natural next question is whether this change in the administratively visible electorate translates into measurable changes in political and policy outcomes.

The answer turns out to be narrower than a standard median-voter story would predict, but not null. Across a broad downstream screen, most outcomes either fail pretrend checks or move only modestly. The patterns that do survive cluster on two margins. Local governments appear to direct more attention toward regular child-oriented schooling, and they expand spending on visible local public goods such as infrastructure and urbanism. National-level partisan vote shares and the broader health-system response do not move in a clean way. I read this configuration as consistent with a single mechanism: BVR did not rewrite local politics, but it changed which residents are easiest for politicians to see and easiest for the local state to reach, and policy responded along exactly those margins.

\subsection{Specification and Screen}
\label{ssec:downstream-strategy}

The downstream estimates follow the baseline design in Section~\ref{sec:empirical-design}: the imputation event-study estimator of \citet{Borusyak2024RevisitiEventStudy} on the strict-BVR sample, with ever-hybrid municipalities excluded, never-treated controls, and municipality-clustered standard errors. The event-time window runs from $-8$ to $+8$, and every reported path is normalized to event time zero so that the post-period values can be read directly from the figures. The screening procedure estimates 298 municipality-level outcomes across political, fiscal, educational, and health domains; 42 pass a pretrend filter, 20 are near, and 236 fail. The main text reports only pass-pretrend outcomes, and the full screen, the ranking tables, and the pretrend-filter rules are described in Appendix~\ref{app:downstream-screen}. Because the screen is broad, I treat the pass-pretrend results as candidates for mechanism evidence rather than as final causal estimates, and I return to the multiple-testing concern at the end of the section. For turnout, blank and null voting, and presidential vote shares, I supplement the imputation paths with dynamic TWFE and Callaway-Sant'Anna estimates as a triangulation, because those outcomes are conceptually closer to the registry contraction and the screen-and-filter logic is less appropriate.

\subsection{Political Outcomes}
\label{ssec:downstream-political}

The cleanest downstream political pattern is mechanical rather than ideological. Turnout, measured as valid votes over registered voters, rises sharply at adoption. The unweighted TWFE estimate at event time zero is about $0.079$, with similar magnitudes at later horizons; the population-weighted estimate is slightly smaller. Because BVR removes voters from the denominator without necessarily changing voting behavior conditional on remaining registered, this is most naturally read as a tighter registry rather than as a behavioral increase in participation. The pattern matches \citet{forquesatoRodrigues2023}, who find that Brazilian biometric recadastramento reduced registered voters but did not significantly change the absolute number of votes cast. Blank-and-null voting does not show a clean post-treatment break, which weakens the alternative reading that voters react to the new technology with confusion or protest.

The more substantive political test asks whether the education-tilted electorate shifted partisan vote shares. The answer is mixed at best. For PT presidential vote share, the unweighted TWFE event study shows negative coefficients at event times $0$ and $4$, the sign one would expect from a less PT-favorable electorate, but the pattern does not survive estimator triangulation. The BJS imputation path is not centered on the same negative effect, the Callaway-Sant'Anna comparison series tells the same near-null story, the population-weighted TWFE estimates are close to zero, and once I absorb state-by-year fixed effects to remove national-cycle shocks, the event-time-zero coefficient flips to positive. The PSDB results are similarly unstable. Figure~\hyperref[fig:downstream-estimator-comparison]{\ref*{fig:downstream-estimator-comparison}} summarizes the estimator comparison for PT vote share alongside one fiscal and one health outcome. I do not read these political patterns as evidence of a partisan realignment. The most credible interpretation is that the registry-composition shift is genuine but too modest, too offsetting across cohorts, or too dominated by national cycles to produce a clean municipality-level partisan downstream effect, consistent with broader evidence that local administration is not a reliable lever on national vote shares \citep{Ferrer2023HowPartisanLocal,Marx2021VoterMobilisaAnd}.

\begin{figure}[htbp]
    \caption{Estimator Comparison for Key Downstream Outcomes}
    \label{fig:downstream-estimator-comparison}
    \centering
    \includegraphics[width=\textwidth]{../resources/images/regressions/estimator_comparison/downstream/downstream_estimator_comparison.pdf}
    \vspace{0.75em}
    \begin{minipage}{\textwidth}
        {\footnotesize \textbf{Note:} The figure overlays the baseline BJS imputation event study \citep{Borusyak2024RevisitiEventStudy} with dynamic TWFE and Callaway-Sant'Anna comparison series for three focal outcomes: PT presidential vote share, log health spending per capita, and infant mortality. Event time is measured relative to first biometric adoption and event time $-2$ is omitted in the TWFE baseline. The comparison series cluster standard errors at the municipality level. \par}
    \end{minipage}
\end{figure}

\newcommand{\BJSFigRoot}{../resources/figures/downstream_bjs_ref0_nonhybrid}

\subsection{Education and Child-Oriented Schooling}
\label{ssec:downstream-education}

The clearest pattern across the downstream screen sits in child-oriented education. Figure~\hyperref[fig:downstream-education]{\ref*{fig:downstream-education}} reports three pass-pretrend BJS event studies. Per-capita spending on child and basic education rises by about $0.09$ log points on average over the post-BVR window, with a flat pre-period and a delayed but positive post-period path. Municipal final-years elementary approval rises by about $1.9$ percentage points, with no detectable pre-trend and a steady upward movement that becomes especially large at longer horizons. Final-years abandonment falls by about $1.0$ percentage point. The screen also reports favorable movement in IDEB approval at the final years of elementary education and in selected SAEB outcomes, although those biennial assessments are noisier; the appendix screen ranking lists them alongside the corresponding pass-pretrend diagnostics.

\begin{figure}[htbp]
    \caption{BVR and Child-Oriented Education}
    \label{fig:downstream-education}
    \centering
    \begin{subfigure}[b]{0.32\textwidth}
        \centering
        \includegraphics[width=\textwidth]{\BJSFigRoot/adult_eja/annual/log_child_spending_pc_bjs_ref0_nonhybrid.pdf}
        \caption{Child/basic spending per capita}
    \end{subfigure}
    \hfill
    \begin{subfigure}[b]{0.32\textwidth}
        \centering
        \includegraphics[width=\textwidth]{\BJSFigRoot/education/annual/inep_ind_municipal_taxa_aprovacao_ef_anos_finais_bjs_ref0_nonhybrid.pdf}
        \caption{Final-years approval}
    \end{subfigure}
    \hfill
    \begin{subfigure}[b]{0.32\textwidth}
        \centering
        \includegraphics[width=\textwidth]{\BJSFigRoot/education/annual/inep_ind_municipal_taxa_abandono_ef_anos_finais_bjs_ref0_nonhybrid.pdf}
        \caption{Final-years abandonment}
    \end{subfigure}
    \vspace{0.5em}
    \begin{minipage}{\textwidth}
        {\footnotesize \textbf{Note:} Each panel reports a pass-pretrend BJS imputation event study normalized to event time zero \citep{Borusyak2024RevisitiEventStudy}. The sample uses strict BVR timing, excludes municipalities ever coded as hybrid, and clusters standard errors at the municipality level. Pre-period diagnostics for each outcome are reported in the appendix screen ranking. \par}
    \end{minipage}
\end{figure}

This pattern is what one would expect if local politicians respond to the post-BVR electorate. The administrative evidence in Sections~\ref{sec:electorate-size} and~\ref{sec:decomposition} shows that the visible electorate became smaller and more educated through the combined channels of low-education exit and education-record updating. Voters whose schooling is updated upward, and voters who complete recadastramento on time, are disproportionately the citizens whose children pass through the regular school system, and they have well-documented preferences over schooling outcomes \citep{Meltzer1981RationalTheoryThe,Hunter2010}. A long tradition in political economy argues that politicians weight the preferences of citizens who are likely to participate and easy to reach \citep{besleyCoate1997,Dixit1996TheDeterminSuccess}. Existing evidence from Brazilian electoral accountability shows that incumbents respond to visible signals about their electorate when those signals are informative \citep{ferrazFinan2011}. The direction here is the opposite of \citet{Fujiwara2015VotingTechnoloPolitica}: electronic voting raised the effective weight of less-educated voters by simplifying the ballot interface, and produced fiscal and health responses favorable to that group; BVR raises the effective weight of more administratively attached voters by tightening pre-election compliance, and the schooling response is consistent with that change in salience.

\subsection{Visible Local Public Goods}
\label{ssec:downstream-public-goods}

The second pass-pretrend pattern is in spending on locally visible public goods. Figure~\hyperref[fig:downstream-public-goods]{\ref*{fig:downstream-public-goods}} reports two SICONFI log per-capita spending outcomes that pass pretrends in the screen: annual infrastructure spending and two-year-cycle urbanism spending. Both rise after adoption, with flat pre-trends and post-period averages of roughly $0.05$ and $0.04$ log points. A related two-year-cycle infrastructure result is near-pass and is reported in the appendix. The corresponding spending \emph{shares} do not pass the pretrend screen for any of these categories, so the evidence is stronger for higher per-capita spending in selected functions than for an active reallocation of the municipal budget toward them.

\begin{figure}[htbp]
    \caption{BVR and Locally Visible Public Goods}
    \label{fig:downstream-public-goods}
    \centering
    \begin{subfigure}[b]{0.48\textwidth}
        \centering
        \includegraphics[width=\textwidth]{\BJSFigRoot/siconfi_log_pc/annual/siconfi_annual_infrastructure_spending_log_pc_bjs_ref0_nonhybrid.pdf}
        \caption{Infrastructure spending per capita}
    \end{subfigure}
    \hfill
    \begin{subfigure}[b]{0.48\textwidth}
        \centering
        \includegraphics[width=\textwidth]{\BJSFigRoot/siconfi_log_pc/cycle2/siconfi_cycle2_urbanism_log_pc_bjs_ref0_nonhybrid.pdf}
        \caption{Urbanism spending per capita}
    \end{subfigure}
    \vspace{0.5em}
    \begin{minipage}{\textwidth}
        {\footnotesize \textbf{Note:} The left panel uses the annual SICONFI spending panel; the right panel uses forward two-year cycles. Both panels show pass-pretrend BJS estimates normalized to event time zero, expressed in $\log(1+x)$ of December 2008 reais per capita. The sample uses strict BVR timing, excludes ever-hybrid municipalities, and clusters standard errors by municipality. \par}
    \end{minipage}
\end{figure}

Two complementary mechanisms fit this pattern, and both connect directly back to the registry results. The first is political: infrastructure and urbanism are visible, locally attributable goods, and a smaller and more educated electorate may reward them more heavily than diffuse transfers \citep{Lindbeck1987BalancedBudgetRedistri,Dixit1996TheDeterminSuccess}. The second is administrative. Biometric recadastramento required municipalities to coordinate with the electoral bureaucracy, set up local service points, and contact citizens. That process plausibly improved the local state's overview of where citizens reside and how they interact with public offices, which makes complementary investments easier to target and to justify. Comparable evidence from biometric monitoring in India shows that authentication and information technologies can improve service delivery and reduce leakage when they also raise the quality of administrative records \citep{Muralidharan2016BuildingStateCapacity,Bossuroy2024BiometriMonitoriService}, and a broader state-capacity literature emphasizes that legibility and effective governance reinforce each other \citep{besleyPersson2011,AcemogluRobinson2019}. The pattern in Figure~\hyperref[fig:downstream-public-goods]{\ref*{fig:downstream-public-goods}} is consistent with both interpretations, and I do not separately identify them here.

\subsection{Birth and Health Outcomes}
\label{ssec:downstream-health}

The health outcomes are more mixed. Figure~\hyperref[fig:downstream-health]{\ref*{fig:downstream-health}} reports three pass- or near-pass outcomes from the screen. BCG coverage rises modestly, the share of low-education mothers with adequate prenatal visits also rises, and the share of births to teen mothers increases by about $1$ percentage point on average over the post window with a flat pre-period and a noticeable positive trajectory at longer horizons. Beyond these, the screen finds favorable movement in low-birthweight outcomes among low-education mothers and in ICU bed availability, but it also reports increases in avoidable mortality and movement in opposite directions for CSAP mortality and SUS beds. The health block as a whole does not look like a broad improvement in service delivery.

\begin{figure}[htbp]
    \caption{Selected Birth and Health Outcomes}
    \label{fig:downstream-health}
    \centering
    \begin{subfigure}[b]{0.32\textwidth}
        \centering
        \includegraphics[width=\textwidth]{\BJSFigRoot/birth_health/annual/sinasc_all_teen_mother_share_known_pp_bjs_ref0_nonhybrid.pdf}
        \caption{Teen mother share}
    \end{subfigure}
    \hfill
    \begin{subfigure}[b]{0.32\textwidth}
        \centering
        \includegraphics[width=\textwidth]{\BJSFigRoot/birth_health/annual/imun_bcg_pp_bjs_ref0_nonhybrid.pdf}
        \caption{BCG coverage}
    \end{subfigure}
    \hfill
    \begin{subfigure}[b]{0.32\textwidth}
        \centering
        \includegraphics[width=\textwidth]{\BJSFigRoot/birth_health/annual/sinasc_lowedu_prenatal_adequate_share_known_pp_bjs_ref0_nonhybrid.pdf}
        \caption{Adequate prenatal, low-ed}
    \end{subfigure}
    \vspace{0.5em}
    \begin{minipage}{\textwidth}
        {\footnotesize \textbf{Note:} BJS event-study paths normalized to event time zero. The teen-mother share is the share of births to mothers under 20 in SINASC. BCG coverage is the percentage of newborns vaccinated. Adequate prenatal visits is the SINASC share of mothers classified as having received an adequate number of prenatal consultations, restricted to mothers with low education. \par}
    \end{minipage}
\end{figure}

The teen-mother result deserves explicit interpretation because it is substantively striking and could be misread. I treat it as most plausibly a measurement-visibility effect of the same administrative process that drives the schooling and infrastructure results. Recadastramento increased local administrative contact with the resident population, and SINASC reporting is more complete in places where the local state interacts more intensively with vulnerable families. A reading in which BVR causally raised teen pregnancy would be hard to reconcile with the absence of broad worsening in other birth or health outcomes and with the favorable movement in low-education prenatal care and BCG coverage. The compositional interpretation, that treated municipalities changed in who gives birth or in whose births are recorded, is harder to rule out and is consistent with the registry composition shift. I report the result here for transparency, but I do not lean on it as a mechanism estimate. Together with the mixed mortality and capacity outcomes, the health evidence is informative mostly as a constraint: BVR did not produce a general health-system response, and any service-delivery response is narrow.

\subsection{A More Legible Electorate}
\label{ssec:downstream-mechanism}

The downstream results fit together as a single story about what BVR changed for politicians and the local state, rather than as a list of separate effects. The registry contraction in Section~\ref{sec:electorate-size} and the count-based decomposition in Section~\ref{sec:decomposition} show that the active electorate became smaller and tilted toward voters recorded as more educated, with positive re-labeling identified in the aggregate and a defensible upper bound on exclusion. Whether the typical low-education registrant disappeared from the rolls or simply had her schooling record updated, the object that politicians and administrators observe in official files is smaller, more selected, and easier to read against the resident population. In a probabilistic-voting setting, the equilibrium policy weights groups that are visible, mobilized, and easier to reach \citep{Lindbeck1987BalancedBudgetRedistri,Dixit1996TheDeterminSuccess}, and turnout-based representation arguments make the same point in a different language \citep{besleyCoate1997,Hajnal2005WhereTurnoutMatters}. The downstream outcomes that move are exactly the ones that mechanism predicts: regular child schooling, which is salient to the post-BVR electorate, and infrastructure and urbanism, which are visible local goods with clear local attribution. The outcomes that do not move are also the ones that mechanism would not predict to respond first: national-cycle-dominated presidential vote shares, broad health-system inputs, and the overall share of the municipal budget. Related Brazilian evidence on restrictions to fraudulent registration shows that comparable shocks to the active electorate can shift candidate selection, public goods, and social outcomes \citep{Karim2025VoterBuyingPolitici}, and the patterns I document are consistent with that broader logic.

\subsection{Interpretation and Limits}
\label{ssec:downstream-limits}

Three caveats discipline the reading of this section. The first is multiple testing. The screen estimates 298 outcomes, and only a small share pass the pretrend filter. Conditioning on a noisy pre-period test induces selection in the post-period contrasts even when each test is individually valid, so the pass-pretrend results should be read as candidate mechanism evidence rather than as final estimates; family-wise and sharpened q-value corrections, alternative pretrend windows, and an explicit pre-registration of a smaller outcome family are reported in Appendix~\ref{app:downstream-screen}. The second caveat is that the level-spending results are stronger than the share results. No SICONFI spending-share outcome passes the pretrend screen in the BJS panel, so the fiscal response is more naturally described as a scale-and-capacity response than as a reallocation of the municipal budget away from other functions. The third caveat is that the LAPOP trust evidence in Section~\ref{sec:political-trust} can also be read through the legibility mechanism: the post-BVR electorate is more selectively composed, and any observed shift in trust among lower-education respondents may partly reflect that selection rather than a within-person belief change. Taken together, the downstream evidence does not argue that BVR transformed Brazilian local government. It argues that BVR changed the part of the electorate that local government most easily observes, and that local policy responded along the margins where that visibility matters most.
